% =============================================================================
\section{Related Work}
\label{sec:related}
% =============================================================================

\subsection{LLM-Based Document Ranking}

LLM-based document ranking has rapidly evolved since the introduction of RankGPT~\citep{sun2023chatgpt}, which demonstrated that properly instructed LLMs can deliver competitive results to supervised rerankers. The field has converged on four prompting paradigms:

\textbf{Pointwise methods}~\citep{nogueira2020document, sachan2022improving} score each document independently. MonoT5~\citep{nogueira2020document} uses a sequence-to-sequence model to predict a ``true'' or ``false'' relevance label for each query-document pair, using the probability of the ``true'' token as a relevance score. While highly efficient---each document requires only one inference call and calls are easily parallelizable---pointwise methods suffer from poor ranking effectiveness because they cannot capture relative relevance between documents. A document that is moderately relevant to one query may receive the same score as one that is highly relevant to another, making cross-document comparison unreliable.

\textbf{Pairwise methods}~\citep{qin2024large, pradeep2021expando} compare two documents at a time, asking the LLM which is more relevant. PRP~\citep{qin2024large} achieves competitive effectiveness with GPT-4-class performance using smaller models like Flan-UL2 (20B), outperforming ChatGPT-based pointwise solutions by 4.2\% on NDCG@10. The DuoT5~\citep{pradeep2021expando} model extends the Expando-Mono-Duo pipeline with a T5-based pairwise reranker. While pairwise methods achieve high effectiveness due to direct inter-document comparison, they require $O(n^2)$ comparisons for all-pairs or $O(n \log n)$ with sorting algorithms---a significant computational overhead.

\textbf{Listwise methods}~\citep{sun2023chatgpt, ma2023zero} present multiple documents and request a complete ranking permutation. RankGPT~\citep{sun2023chatgpt} uses a sliding window strategy: a window of $w$ documents slides from the bottom of the initial ranking to the top, with the LLM reranking each window. This approach captures rich inter-document relationships and produces coherent rankings. However, it is constrained by context length, sensitive to input order (position bias), and expensive due to long prompts. \citet{zhang2025rankwithoutgpt} demonstrate that effective listwise rerankers can be built on open-source LLMs, achieving 97\% of GPT-4 effectiveness.

\textbf{Setwise methods}~\citep{zhuang2024setwise} ask the LLM to select the most relevant document from a small candidate set (typically 3--4 documents), using the selection as a comparator in sorting algorithms. This paradigm balances effectiveness and efficiency: it extracts more comparative information per call than pairwise (multiple documents compared simultaneously) while using shorter prompts than listwise (only a small subset presented). The approach has been extended by setwise insertion~\citep{podolak2025beyond}, which leverages the initial BM25 ranking as prior knowledge for 31\% time reduction; tournament-based ranking (TourRank)~\citep{chen2025tourrank}, which uses sports tournament brackets for more robust rankings; and reasoning-enhanced variants (Rank-R1)~\citep{zhuang2025rankr1}, which integrate chain-of-thought reasoning via reinforcement learning.

All existing setwise methods share a common design choice: they only ask the LLM to identify the \emph{most relevant} document. Our work is the first to investigate alternative selection targets (least relevant) and simultaneous multi-target selection (both most and least), along with the sorting algorithms that consume those richer outputs.

\subsection{Sorting Algorithms for LLM Ranking}

Classical sorting algorithms have been adapted as the backbone for pairwise and setwise LLM ranking. The choice of algorithm significantly impacts both the number of LLM calls and the quality of the final ranking:

\textbf{Heapsort} builds an $n$-ary max-heap where each node is compared against its $c$ children (for setwise with \texttt{num\_child}$=c$). Building the heap requires $O(n)$ comparisons, and extracting $k$ maximums requires $O(k \log_c n)$ comparisons, yielding $O(n + k \log_c n)$ total. For the typical setting ($k=10$, $n=100$, $c=3$), this is approximately 130 comparisons.

\textbf{Bubblesort} starts from the bottom of the list and ``bubbles'' the most relevant document to the top through a series of windowed comparisons. Each of $k$ passes requires $O(n/c)$ comparisons, giving $O(kn/c)$ total---approximately 333 comparisons for the same setting. However, an optimization tracking the last-start position can reduce this significantly by skipping unchanged regions.

Recent work has challenged whether classical complexity bounds remain optimal when the comparator is an LLM. \citet{sato2026survey} surveys sorting algorithms adapted for LLMs, noting that traditional cost models fail to account for three key properties: (1) LLM comparisons are expensive in both time and cost, (2) comparisons can be batched for parallel execution, and (3) comparisons are noisy and potentially non-transitive. \citet{wisznia2025optimal} further show that batching and caching fundamentally change the cost model, making some ``suboptimal'' algorithms superior in practice.

Tournament-based approaches like TourRank~\citep{chen2025tourrank} and BlitzRank~\citep{blitzrank2026} represent a move toward extracting richer information from each comparison round. BlitzRank models the complete preference graph from $k$-wise comparisons rather than discarding loser information. However, these approaches still focus on identifying winners. Our cocktail shaker sort and double-ended selection sort are the first sorting algorithms specifically designed to exploit dual-output (best \emph{and} worst) comparators.

\subsection{Position Bias in LLM Ranking}

LLMs exhibit systematic bias based on document position in the prompt. The ``Lost in the Middle'' phenomenon~\citep{liu2024lost} shows that LLMs attend more to information at the beginning and end of long contexts, neglecting middle positions. For ranking, this translates to position-dependent selection probabilities that can distort the final ranking.

Several mitigation strategies have been proposed. Permutation self-consistency~\citep{tang2024found} shuffles document order across multiple passes and aggregates rankings via voting, achieving 7--18\% improvement for GPT-3.5 and 8--16\% for LLaMA v2 (70B). However, this multiplies computational cost by the number of permutations. LLM-RankFusion~\citep{zeng2024llmrankfusion} identifies two types of intrinsic inconsistency: \emph{order inconsistency} (conflicting results when document order is swapped) and \emph{transitive inconsistency} (non-transitive preference triads where $A > B$, $B > C$, but $C > A$). These inconsistencies fundamentally challenge sorting-based approaches, which assume a consistent total order.

For setwise ranking specifically, position bias manifests within the small candidate set. \citet{zhuang2024setwise} mitigate this through permutation voting: randomly shuffling both document order and label assignment across multiple permutations, with majority voting to determine the winner.

Our work opens a new dimension of analysis: whether position bias operates differently for best-selection versus worst-selection prompts, and---more interestingly---whether the bias for the \emph{same} model on the \emph{same} candidate set changes when ``best'' and ``worst'' are asked jointly versus alone. As we show in \S\ref{sec:bias}, the dual prompt produces a reversed dual\_worst pattern that is consistent across all 9 of our models and gives an empirically grounded mechanism for the directional asymmetry.

\subsection{Negative Evidence in LLM Ranking}

A growing body of LLM-ranking work also reports negative or surprising findings. \citet{zeng2024llmrankfusion} document order- and transitivity-inconsistencies. \citet{tang2024found} treat position bias as a defect and use permutation voting to suppress it. We add to this literature with two coherent negative results: (i) standalone \bottomup{} is consistently weaker than \topdown{}, and (ii) directional fusion (\bidir{}) does not improve over the better individual ranker. These are not isolated tuning failures: they reproduce across 9 models, 2 datasets, multiple fusion methods, and the full range of weight $\alpha$, and they survive Bonferroni correction in 6/18 and 3/18 configurations respectively (\S\ref{sec:rq1}, \S\ref{sec:rq3}).
